Krátký praktický checklist — rychlá kontrola před i během použití generativní AI ve výuce.

\begin{itemize}
\item \textbf{Před nasazením:} Proveďte DPIA pokud se zpracovávají studentská data a uveďte ji při vyjednávání DPA;\cite{kb_ai_101_tipu_pdf_04e4a115_page_8_1} zkontrolujte DPA (role, sub‑procesory, retention, audit) a vyžadujte SLA pro incident response;\cite{kb_ai_101_tipu_pdf_04e4a115_page_8_1} upřednostněte enterprise/licencované varianty s možností hostingu v přijatelné jurisdikci;\cite{kb_ai_101_tipu_pdf_04e4a115_page_3_1} minimalizujte sběr osobních údajů a neodesílejte studentské fotografie do neznámých cloudů.\cite{kb_ai_101_tipu_pdf_04e4a115_page_53_2}
\item \textbf{Během provozu:} Pseudonymizujte nebo agregujte data, nastavte role a přístupy s loggingem/monitoringem a sledujte používání a chybové stavy;\cite{kb_ai_101_tipu_pdf_04e4a115_page_8_1} mějte připravený postup incident response (kontakty IT/právo, notifikace podle DPA/SLA).\cite{kb_ai_101_tipu_pdf_04e4a115_page_8_1}
\item \textbf{Krátký „cheat‑sheet“ pro incident:} Detect, Contain, Notify, Remediate.
\end{itemize}